% Code listing configuration for MATLAB
\lstset{
    language=Matlab,
    basicstyle=\ttfamily\small,
    keywordstyle=\color{blue},
    stringstyle=\color{purple},
    commentstyle=\color{green!60!black},
    numbers=left,
    numberstyle=\tiny\color{gray},
    breaklines=true,
    frame=single,
    captionpos=b
}

\section{Exercise 7.10: DFT basics}
\subsection{MATLAB Implementation}

Below is the MATLAB script dft\_delta.m. This script performs a Discrete Fourier Transform (DFT) calculation on a simple array labeled f consisting of zeros with a length of 10 and the first value set to one. The script then computes the Inverse Discrete Fourier Transform (IDFT) on the the subsequent array labeled F and achieves the original array f. 

\begin{lstlisting}[caption={MATLAB script dft_delta.m for DFT basics}, label={lst:dft_delta}]

% 1. Declare variable N and array f. Set f(1) to 1, all other entries 0.
N = 10;
f = zeros(1, N); f(1) = 1; 

% 2. Calculate the DFT of f and print it.
F = fft(f);
F

% 3. Calculate the IDFT of F and print it.
F_inv = ifft(F);
F_inv

% 4. Express every Omega_k terms of k and N. 
k = 0:N-1;
Omega_k = 2*pi*k / N;

\end{lstlisting}

The DFT of a delta-peak function is F[k] = 1, here we see that same result when performing a DFT on f.